\subsection{Examples of locally convex spaces}

\begin{enumerate}
\def\labelenumi{\arabic{enumi})}
\item
  The space \(R^{n}\) is locally convex since the open cubes with centre 0 are convex (II, p.~9, prop. 6). This is, therefore, also true for all real topological vector spaces of finite dimension; in fact it follows from the above and I, § 2.3, th. 2 provided that E is Hausdorff; if not, the Hausdorff space F associated with E is of finite dimension, therefore locally convex, and the inverse images of convex neighbourhoods of 0 in F under the canonical mapping \(E \rightarrow F\) are convex and form a fundamental system of neighbourhoods of 0 in E.
\item
  Let E be a vector space in R, and B be the family of all subsets of E that are absorbent, symmetric and convex. By prop. 1 of II, p.~23 we see that B is a fundamental system of neighbourhoods of 0 for a locally convex topology \(T_{\omega}\) on E that is the finest of all locally convex topologies on E. This topology is Hausdorff; for let \(x \neq 0\) be any point of E; there exists a basis \((i_{\mathfrak{t}})_{\mathfrak{t} \in \mathbf{I}}\) of E with an \(\alpha \in \mathbf{I}\) such that \(e_{\alpha} = x\) ; the set of points \(y = \sum_{\mathfrak{t}} y_{\mathfrak{t}} e_{\mathfrak{t}}\) such that \(|y_{\alpha}| < 1\) is absorbent, symmetric and convex. It does not contain x.
\end{enumerate}

From II, p.~24, corollary, it follows that \(T_{\omega}\) is also the topology defined by the set of all semi-norms on E, thus every semi-norm is continuous in \(T_{\omega}\) .

In particular, if u is a linear mapping of E in any locally convex space F, the inverse image, under u, of every convex neighbourhood of 0 in F is an absorbent convex set in E; therefore it is a neighbourhood of 0 for \(T_{\omega}\) and thus u is continuous for \(T_{\omega}\) .

Given a convex set C in E, we say that a point \(a \in C\) is an internal point of C if, for every line D containing a, the intersection \(D \cap C\) contains an open segment which contains a; in other words \(-a + C\) is absorbent. The point a of the set A in E is interior to A for \(T_{\omega}\) if, and only if, there exists a convex set C with \(a \in C \subset A\) , and such that a is an internal point of C.

More generally, let V be an affine linear variety in E, and C be a convex set contained in V; a point \(a \in C\) is an internal point of C relative to V if, in the vector subspace \(V_{0} = -a + V\) , the point 0 is an internal point of the set \(C_{0} = -a + C\) .

When E is of finite dimension, the topology \(T_{\omega}\) is just the canonical topology on E (I, p.~13, th. 2); which shows that every internal point of a convex set C in E, is interior to C for the canonical topology (cf.~II, p.~74, exerc. 5).

\begin{enumerate}
\def\labelenumi{\arabic{enumi})}
\setcounter{enumi}{2}
\tightlist
\item
  Let A be a symmetric convex set in the vector space E over R. The vector subspace F generated by A is also the convex cone generated by A, since -A = A; this set is the set of \(\lambda x\) where \(x \in A\) and \(\lambda \in R\) ; the set A is absorbent in F and the sets \(\lambda A\) where \(\lambda > 0\) , form a fundamental system of neighbourhoods of 0 for a locally convex topology on F (said to be defined by A), which is defined by the semi-norm \(p_{A}\) , the gauge of A (II, p.~20, prop. 22); we write \(E_{A}\) for the locally convex space obtained by giving F this semi-norm. The space \(E_{A}\) is Hausdorff if, and only if, \(p_{A}\) is a norm or alternatively A does not contain any line. If B is a second symmetric convex set in E and if \(A \subset B\) , then clearly \(E_{A} \subset E_{B}\) , and the canonical injection of \(E_{A}\) in \(E_{B}\) is continuous for the topologies defined respectively by A and by B. Further, if f is a linear mapping of E in a real vector space \(E'\) , then \(f(A)\) is convex and symmetric in \(E'\) and f is a continuous linear mapping of \(E_{A}\) on \(E'_{f(A)}\) .
\end{enumerate}

Finally, note that if E carries a topology T compatible with its vector space structure, and if V is a symmetric convex neighbourhood of 0 for T, then the vector space generated by V is identical with E, since V is absorbent, and the identity mapping of E in \(E_{V}\) is continuous.

